\documentclass[11pt]{article}
\usepackage[utf8]{inputenc}
\usepackage{ctex}
\usepackage{amsmath, amssymb, amsthm}
\usepackage{geometry}
\geometry{a4paper, margin=1in}

\input{../preamable/preamable_sep.tex}

\declaretheorem[style=thmgreenbox, name=定义]{cndefinition}
\declaretheorem[style=thmredbox, name=定理]{cntheorem}
\declaretheorem[style=thmbluebox, name=性质]{cnproperty}
\declaretheorem[style=thmredbox, name=引理]{cnlemma}
\declaretheorem[style=thmredbox, numbered=no, name=推论]{cncorollary}
\declaretheorem[style=thmbluebox, numbered=no, name=例]{cnexample}
\title{近世代数笔记(Lec 15)：有限域}
\author{Raysance}
\date{}

\begin{document}

\maketitle

\section{有限域的基本性质}

有限域（Finite Field），又称伽罗瓦域（Galois Field），是代数学中研究极其广泛的对象。在这一节中，我们首先探讨有限域的一些基础构造与代数性质。

\begin{cnproperty}
    若 $F$ 是一个有限域，则其特征 $\text{char}(F) = p$ 必为一个素数。
\end{cnproperty}
\begin{proof}
    由域的性质可知，有限域的单位元 $1$ 的阶数必然有限（否则 $F$ 包含一个与 $\mathbb{Z}$ 同构的子环，从而无限）。设 $\text{char}(F) = m$。如果 $m = ab$ 是合数（$a, b > 1$），则 $(a \cdot 1)(b \cdot 1) = ab \cdot 1 = 0$。由于域中无零因子，必有 $a \cdot 1 = 0$ 或 $b \cdot 1 = 0$，这与 $m$ 是特征的定义（最小正整数）矛盾。因此 $m$ 必为素数 $p$。
\end{proof}

\begin{cntheorem}
    设 $F$ 是有限域，且 $\text{char}(F) = p$，则存在正整数 $n$ 使得 $|F| = p^n$。
\end{cntheorem}
\begin{proof}
    由于 $\text{char}(F) = p$，$F$ 必包含一个同构于 $\mathbb{Z}_p$ 的子域（素域）。此时 $F$ 可以看作是 $\mathbb{Z}_p$ 上的向量空间。因为 $F$ 有限，其维数 $n = [F : \mathbb{Z}_p]$ 也必然有限。设 $\{v_1, \dots, v_n\}$ 是 $F$ 作为 $\mathbb{Z}_p$-向量空间的一组基，则 $F$ 中的任意元素都可以唯一表示为：
    \[ x = a_1 v_1 + a_2 v_2 + \dots + a_n v_n, \quad a_i \in \mathbb{Z}_p \]
    由于每个系数 $a_i$ 都有 $p$ 种选择，故 $|F| = p^n$。
\end{proof}

\begin{cntheorem}
    有限域 $F$ 的乘法群 $F^*$ 是循环群。
\end{cntheorem}
\begin{proof}
    设 $|F| = q = p^n$，则 $|F^*| = q-1$。根据有限阿贝尔群的结构定理，$F^*$ 同构于若干个循环群的直和：
    \[ F^* \cong C_{d_1} \times C_{d_2} \times \dots \times C_{d_k} \]
    其中 $d_1 | d_2 | \dots | d_k$。由此可知对于任意 $x \in F^*$，都有 $x^{d_k} = 1$。这意味着多项式 $f(x) = x^{d_k} - 1$ 在域 $F$ 中有 $|F^*| = q-1$ 个根。然而，次数为 $d_k$ 的多项式在域中至多有 $d_k$ 个根，因此 $q-1 \le d_k$。又因为 $d_k$ 是 $F^*$ 的子群阶数，必有 $d_k \le q-1$。综上所述 $d_k = q-1$。这意味着 $F^*$ 中存在阶为 $q-1$ 的元素，即 $F^*$ 是循环群。
\end{proof}

\section{有限域的存在性与唯一性}

接下来的核心问题是：给定 $q = p^n$，是否一定存在一个阶为 $q$ 的域？这个域是否唯一？

\begin{cntheorem}
    对于任意素数 $p$ 和正整数 $n$，在 $\mathbb{Z}_p$ 的代数闭包中，恰好存在一个 $p^n$ 阶的子域。它是多项式 $f(x) = x^{p^n} - x$ 在 $\mathbb{Z}_p$ 上的分裂域。
\end{cntheorem}

\begin{proof}
    \textbf{1. 无重根性}：考虑多项式 $f(x) = x^q - x$，其中 $q = p^n$。其导数为 $f'(x) = qx^{q-1} - 1 = n p x^{q-1} - 1$。由于 $\text{char}(F) = p$，在 $\mathbb{Z}_p$ 中 $p = 0$，故 $f'(x) = -1$。由于 $f'(x)$ 恒不为零，它与 $f(x)$ 显然互素，即 $\gcd(f(x), f'(x)) = 1$。这说明 $x^q - x$ 在其分裂域中没有重根，共有 $q$ 个不同的根。

    \textbf{2. 根集构成域}：令 $S$ 为 $x^q - x$ 在其分裂域中的所有根构成的集合。对于任意 $a, b \in S$，有 $a^q = a$ 且 $b^q = b$。则：
    \begin{itemize}
        \item $(a \pm b)^q = a^q \pm b^q = a \pm b \implies a \pm b \in S$（利用 Frobenius 映射的线性性质）；
        \item $(ab)^q = a^q b^q = ab \implies ab \in S$；
        \item 对 $b \neq 0$，$(a/b)^q = a^q / b^q = a/b \implies a/b \in S$。
    \end{itemize}
    因此 $S$ 构成一个域。由于 $S$ 包含 $q$ 个元素且包含 $\mathbb{Z}_p$（因为 $\forall a \in \mathbb{Z}_p, a^p = a \implies a^{p^n} = a$），它就是 $x^q - x$ 在 $\mathbb{Z}_p$ 上的分裂域。

    \textbf{3. 唯一性}：根据分裂域在同构意义下的唯一性，任何 $q$ 阶域必然是 $x^q - x$ 的分裂域，因此它们也相互同构。我们将这个阶为 $q$ 的有限域记作 $GF(q)$ 或 $F_q$。
\end{proof}

\section{有限域的子域结构}

有限域之间的包含关系与其阶数有着紧密的逻辑联系。

\begin{cntheorem}
    设 $q = p^k$，则 $F_q$ 是 $F_{q^n}$ 的子域当且仅当 $k$ 整除 $n$。
\end{cntheorem}
\begin{proof}
    该命题可以转化为多项式的整除问题。
    首先，$F_{p^k} \subseteq F_{p^n}$ 的充要条件是 $x^{p^k} - x$ 的所有根都是 $x^{p^n} - x$ 的根，即 $x^{p^k} - x \mid x^{p^n} - x$。
    
    引理：在 $\mathbb{Z}$ 中，$a^k - 1 \mid a^n - 1 \iff k \mid n$。
    证明：设 $n = mk + r$ ($0 \le r < k$)。则 $a^n - 1 = a^{mk+r} - 1 = a^r(a^{mk} - 1) + (a^r - 1)$。由于 $a^k - 1 \mid (a^k)^m - 1$，故 $a^k - 1 \mid a^n - 1 \iff a^k - 1 \mid a^r - 1$。因为 $r < k$，这只有在 $r=0$ 时成立。

    应用此引理：
    $x^{p^k-1} - 1 \mid x^{p^n-1} - 1 \iff p^k - 1 \mid p^n - 1 \iff k \mid n$。
    而 $x^{p^k} - x \mid x^{p^n} - x \iff x(x^{p^k-1} - 1) \mid x(x^{p^n-1} - 1) \iff x^{p^k-1} - 1 \mid x^{p^n-1} - 1$。
    因此，$F_{p^k} \subseteq F_{p^n} \iff k \mid n$。
\end{proof}

\section{有限域的 Galois 理论}

有限域上的代数扩张是非常“漂亮”的，因为它们总是 Galois 扩张，且 Galois 群具有极其简单的循环结构。

\begin{cntheorem}
    扩张 $F_{q^n} / F_q$ 是 Galois 扩张，且其 Galois 群 $\text{Gal}(F_{q^n}/F_q)$ 是一个 $n$ 阶循环群。
\end{cntheorem}
\begin{proof}
    \textbf{1. 扩张性质}：$F_{q^n}$ 是 $x^{q^n} - x$ 在 $F_q$ 上的分裂域（因为 $F_q$ 是 $F_{q^n}$ 的子域），且该多项式无重根。因此该扩张是正规且可分的，即为 Galois 扩张。

    \textbf{2. Frobenius 自同构}：考虑映射 $\sigma : F_{q^n} \to F_{q^n}, \alpha \mapsto \alpha^q$。
    \begin{itemize}
        \item $\sigma$ 是域自同构：$\sigma(a+b) = (a+b)^q = a^q + b^q$（特征为 $p$ 且 $q$ 是 $p$ 的幂），$\sigma(ab) = (ab)^q = a^q b^q$。
        \item $\sigma$ 保持 $F_q$ 不变：$\forall a \in F_q, a^q = a \implies \sigma(a) = a$。
    \end{itemize}
    因此 $\sigma \in \text{Gal}(F_{q^n}/F_q)$。

    \textbf{3. 群的结构}：$\sigma^k(\alpha) = \alpha^{q^k}$。如果 $\sigma^k = \text{id}$，则对所有 $\alpha \in F_{q^n}$ 都有 $\alpha^{q^k} - \alpha = 0$。这意味着 $x^{q^k} - x$ 有 $q^n$ 个根，从而 $q^k \ge q^n$，即 $k \ge n$。因此 $\sigma$ 的阶至少为 $n$。又因为 $|\text{Gal}(F_{q^n}/F_q)| = [F_{q^n} : F_q] = n$，故 $\sigma$ 是该群的生成元，该群是循环群。我们称 $\sigma$ 为 \textbf{Frobenius 自同构}。
\end{proof}

\section{有限域的算术与例子：$F_{16}$}

在实际计算中，我们常通过不可约多项式来构造有限域。

\subsection{极小多项式与共轭元}
设 $\alpha \in F_{q^n}$。与 $\alpha$ 关于 $F_q$ 共轭的元素是该元素在 Frobenius 映射下的像：
\[ \alpha, \alpha^q, \alpha^{q^2}, \dots, \alpha^{q^{d-1}} \]
其中 $d$ 是使得 $\alpha^{q^d} = \alpha$ 的最小正整数，它等于 $\alpha$ 在 $F_q$ 上的极小多项式的次数，且 $d \mid n$。
$\alpha$ 在 $F_q$ 上的极小多项式为：
\[ m_\alpha(x) = (x - \alpha)(x - \alpha^q) \dots (x - \alpha^{q^{d-1}}) \]

\subsection{实例：$F_{16} / F_2$}
我们利用 $F_2$ 上的 4 次不可约多项式 $p(x) = x^4 + x + 1$ 来构造 $F_{16}$。
令 $\alpha$ 是 $p(x)$ 的一个根，则 $\alpha^4 = \alpha + 1$。$F_{16}$ 中的元素可以表示为 $a_3 \alpha^3 + a_2 \alpha^2 + a_1 \alpha + a_0$。

由于 $F_{16}^*$ 是 15 阶循环群，$\alpha$ 是其生成元（本原元）。我们可以列出幂表：

\begin{table}[h]
\centering
\begin{tabular}{|l|l|l|}
\hline
指数 $k$ & 向量表示 $(a_3 a_2 a_1 a_0)$ & 多项式表示 \\ \hline
0 (1)      & 0001 & 1 \\ \hline
1 ($\alpha$) & 0010 & $\alpha$ \\ \hline
2 ($\alpha^2$) & 0100 & $\alpha^2$ \\ \hline
3 ($\alpha^3$) & 1000 & $\alpha^3$ \\ \hline
4 ($\alpha^4$) & 0011 & $\alpha + 1$ \\ \hline
5 ($\alpha^5$) & 0110 & $\alpha^2 + \alpha$ \\ \hline
6 ($\alpha^6$) & 1100 & $\alpha^3 + \alpha^2$ \\ \hline
7 ($\alpha^7$) & 1011 & $\alpha^3 + \alpha + 1$ \\ \hline
8 ($\alpha^8$) & 0101 & $\alpha^2 + 1$ \\ \hline
9 ($\alpha^9$) & 1010 & $\alpha^3 + \alpha$ \\ \hline
10 ($\alpha^{10}$) & 0111 & $\alpha^2 + \alpha + 1$ \\ \hline
11 ($\alpha^{11}$) & 1110 & $\alpha^3 + \alpha^2 + \alpha$ \\ \hline
12 ($\alpha^{12}$) & 1111 & $\alpha^3 + \alpha^2 + \alpha + 1$ \\ \hline
13 ($\alpha^{13}$) & 1101 & $\alpha^3 + \alpha^2 + 1$ \\ \hline
14 ($\alpha^{14}$) & 1001 & $\alpha^3 + 1$ \\ \hline
\end{tabular}
\end{table}

\paragraph{子域与极小多项式分解}
$x^{15} - 1$ 在 $F_2$ 上的分解对应于各元素的极小多项式乘积。
\begin{itemize}
    \item $\alpha$ 的共轭元：$\alpha, \alpha^2, \alpha^4, \alpha^8$。极小多项式 $m_1(x) = x^4 + x + 1$。
    \item $\alpha^3$ 的共轭元：$\alpha^3, \alpha^6, \alpha^{12}, \alpha^9$。极小多项式 $m_3(x) = (x-\alpha^3)(x-\alpha^6)(x-\alpha^{12})(x-\alpha^9) = x^4 + x^3 + x^2 + x + 1$。
    \item $\alpha^5$ 的共轭元：$\alpha^5, \alpha^{10}$。极小多项式 $m_5(x) = (x-\alpha^5)(x-\alpha^{10}) = x^2 + x + 1$（对应子域 $F_4$）。
    \item $\alpha^7$ 的共轭元：$\alpha^7, \alpha^{14}, \alpha^{28}=\alpha^{13}, \alpha^{26}=\alpha^{11}$。极小多项式 $m_7(x) = x^4 + x^3 + 1$。
    \item $1$ 的极小多项式：$m_0(x) = x + 1$。
\end{itemize}
综上，$x^{15}-1 = m_0(x)m_1(x)m_3(x)m_5(x)m_7(x)$。

\end{document}
